\documentclass[10pt,a4paper]{amsart}
\usepackage{amsmath,amssymb,amsthm,graphicx}
% Retained preamble for the Stacks Project chapter {CH} (shared project preamble).
% The chapter file pinned-original-source.tex only carries \input{{preamble}} on its first line,
% so no preamble text lives inside the chapter itself.  The Stacks Project uses one single shared
% preamble for every chapter: the snapshot stages it three times (sources/stacks-chow/preamble.tex,
% sources/stacks-etale-cohomology/preamble.tex and focused-seeds/AGP/sources/stacks-snapshot-20261011/
% preamble.tex) and those copies are identical in size.  The block below is copied mechanically,
% line for line, from sources/stacks-chow/preamble.tex lines 14-18 and 154-217: the Xy-pic setup and
% the native theorem declarations together with the mathematical macros.  Nothing was added, removed
% or redefined; the wrapper class and the amsmath/amssymb/amsthm/graphicx base packages are supplied
% by the assembler.
% For commutative diagrams we use Xy-pic
% Chapter: stacks-spaces-more-morphisms.tex -- More on Morphisms of Spaces.
\usepackage[all]{xy}

% We use 2cell for 2-commutative diagrams.
\xyoption{2cell}

% Theorem environments.
%
\theoremstyle{plain}
\newtheorem{theorem}[subsection]{Theorem}
\newtheorem{proposition}[subsection]{Proposition}
\newtheorem{lemma}[subsection]{Lemma}

\theoremstyle{definition}
\newtheorem{definition}[subsection]{Definition}
\newtheorem{example}[subsection]{Example}
\newtheorem{exercise}[subsection]{Exercise}
\newtheorem{situation}[subsection]{Situation}

\theoremstyle{remark}
\newtheorem{remark}[subsection]{Remark}
\newtheorem{remarks}[subsection]{Remarks}

\numberwithin{equation}{subsection}

% Macros
%
\def\lim{\mathop{\mathrm{lim}}\nolimits}
\def\colim{\mathop{\mathrm{colim}}\nolimits}
\def\Spec{\mathop{\mathrm{Spec}}}
\def\Hom{\mathop{\mathrm{Hom}}\nolimits}
\def\Ext{\mathop{\mathrm{Ext}}\nolimits}
\def\SheafHom{\mathop{\mathcal{H}\!\mathit{om}}\nolimits}
\def\SheafExt{\mathop{\mathcal{E}\!\mathit{xt}}\nolimits}
\def\Sch{\mathit{Sch}}
\def\Mor{\mathop{\mathrm{Mor}}\nolimits}
\def\Ob{\mathop{\mathrm{Ob}}\nolimits}
\def\Sh{\mathop{\mathit{Sh}}\nolimits}
\def\NL{\mathop{N\!L}\nolimits}
\def\CH{\mathop{\mathrm{CH}}\nolimits}
\def\proetale{{pro\text{-}\acute{e}tale}}
\def\etale{{\acute{e}tale}}
\def\QCoh{\mathit{QCoh}}
\def\Ker{\mathop{\mathrm{Ker}}}
\def\Im{\mathop{\mathrm{Im}}}
\def\Coker{\mathop{\mathrm{Coker}}}
\def\Coim{\mathop{\mathrm{Coim}}}

% Boxtimes
%
\DeclareMathSymbol{\boxtimes}{\mathbin}{AMSa}{"02}

%
% Macros for moduli stacks/spaces
%
\def\QCohstack{\mathcal{QC}\!\mathit{oh}}
\def\Cohstack{\mathcal{C}\!\mathit{oh}}
\def\Spacesstack{\mathcal{S}\!\mathit{paces}}
\def\Quotfunctor{\mathrm{Quot}}
\def\Hilbfunctor{\mathrm{Hilb}}
\def\Curvesstack{\mathcal{C}\!\mathit{urves}}
\def\Polarizedstack{\mathcal{P}\!\mathit{olarized}}
\def\Complexesstack{\mathcal{C}\!\mathit{omplexes}}
% \Pic is the operator that assigns to X its picard group, usage \Pic(X)
% \Picardstack_{X/B} denotes the Picard stack of X over B
% \Picardfunctor_{X/B} denotes the Picard functor of X over B
\def\Pic{\mathop{\mathrm{Pic}}\nolimits}
\def\Picardstack{\mathcal{P}\!\mathit{ic}}
\def\Picardfunctor{\mathrm{Pic}}
\def\Deformationcategory{\mathcal{D}\!\mathit{ef}}
\usepackage{cleveref}
\crefname{originallemma}{original lemma}{original lemmas}
\Crefname{originallemma}{Original lemma}{Original lemmas}
\makeatletter
\newlabel{algebra-lemma-quotient-strict-henselization}{{algebra-lemma-quotient-strict-henselization}{source}{}{sourcepointer-1}{}}
\newlabel{algebra-lemma-quotient-strict-henselization@cref}{{[originallemma][0][]algebra-lemma-quotient-strict-henselization}{[1][0][]source}{}{}{}}
\makeatother
\crefname{originallemma}{original lemma}{original lemmas}
\Crefname{originallemma}{Original lemma}{Original lemmas}
\makeatletter
\newlabel{spaces-morphisms-definition-flat-module}{{spaces-morphisms-definition-flat-module}{source}{}{sourcepointer-2}{}}
\newlabel{spaces-morphisms-definition-flat-module@cref}{{[originallemma][0][]spaces-morphisms-definition-flat-module}{[1][0][]source}{}{}{}}
\makeatother
\crefname{originallemma}{original lemma}{original lemmas}
\Crefname{originallemma}{Original lemma}{Original lemmas}
\makeatletter
\newlabel{spaces-morphisms-lemma-base-change-module-flat}{{spaces-morphisms-lemma-base-change-module-flat}{source}{}{sourcepointer-3}{}}
\newlabel{spaces-morphisms-lemma-base-change-module-flat@cref}{{[originallemma][0][]spaces-morphisms-lemma-base-change-module-flat}{[1][0][]source}{}{}{}}
\makeatother
\crefname{originallemma}{original lemma}{original lemmas}
\Crefname{originallemma}{Original lemma}{Original lemmas}
\makeatletter
\newlabel{spaces-morphisms-lemma-flat-at-point}{{spaces-morphisms-lemma-flat-at-point}{source}{}{sourcepointer-4}{}}
\newlabel{spaces-morphisms-lemma-flat-at-point@cref}{{[originallemma][0][]spaces-morphisms-lemma-flat-at-point}{[1][0][]source}{}{}{}}
\makeatother
\crefname{originallemma}{original lemma}{original lemmas}
\Crefname{originallemma}{Original lemma}{Original lemmas}
\makeatletter
\newlabel{spaces-properties-lemma-describe-etale-local-ring}{{spaces-properties-lemma-describe-etale-local-ring}{source}{}{sourcepointer-5}{}}
\newlabel{spaces-properties-lemma-describe-etale-local-ring@cref}{{[originallemma][0][]spaces-properties-lemma-describe-etale-local-ring}{[1][0][]source}{}{}{}}
\makeatother
\crefname{originallemma}{original lemma}{original lemmas}
\Crefname{originallemma}{Original lemma}{Original lemmas}
\makeatletter
\newlabel{spaces-properties-lemma-geometric-lift-to-usual}{{spaces-properties-lemma-geometric-lift-to-usual}{source}{}{sourcepointer-6}{}}
\newlabel{spaces-properties-lemma-geometric-lift-to-usual@cref}{{[originallemma][0][]spaces-properties-lemma-geometric-lift-to-usual}{[1][0][]source}{}{}{}}
\makeatother
\crefname{originallemma}{original lemma}{original lemmas}
\Crefname{originallemma}{Original lemma}{Original lemmas}
\makeatletter
\newlabel{spaces-properties-lemma-stalk-quasi-coherent}{{spaces-properties-lemma-stalk-quasi-coherent}{source}{}{sourcepointer-7}{}}
\newlabel{spaces-properties-lemma-stalk-quasi-coherent@cref}{{[originallemma][0][]spaces-properties-lemma-stalk-quasi-coherent}{[1][0][]source}{}{}{}}
\makeatother
\title{Three equivalent formulations of flatness of a sheaf on the fibres of a morphism of algebraic spaces}
\author{The Stacks Project Authors}
\date{stacks-project:spaces-more-morphisms.tex:local-snapshot-20261011; GNU FDL 1.2 or later}
\begin{document}
\maketitle
\noindent\textit{Editorial note.} Statement span is lines 5631-5663, including the large xymatrix display inside item (3); proof span is lines 5665-5787.\par\medskip
\noindent\textit{Editorial note.} The statement's xymatrix uses an unusual arrow specification; it is retained verbatim and the retained preamble provides the Xy-pic setup.\par\medskip
\noindent\textit{Editorial note.} The proof opens with a parenthetical 'details omitted' for the independence of condition (2) of the chosen point, and then says this will also follow from the arguments below; both are retained verbatim. The main comparison of the three conditions is carried out in full.\par\medskip
\noindent\textit{Editorial note.} The proof cites several labelled results, which are prerequisites and are supplied as non-load-bearing navigation stubs.\par\medskip
\noindent\textit{Editorial note.} extra\_packages is empty on purpose: the retained preamble already loads xy with options and a second declaration clashes.\par\medskip
\noindent\textit{Allowed standard prerequisites (editorial).}\begin{itemize}
\item Explicit prerequisite boundary. Etale morphisms of algebraic spaces and surjective etale charts, geometric points and their liftings, strict henselisation of a local ring, and the definition of flatness for a module on an algebraic space. The comparison of the three conditions is retained.
\end{itemize}\medskip
\ifdefined\hypertarget\hypertarget{sourcepointer-1}{}\fi
\noindent\textit{Source reference.} Original lemma algebra-lemma-quotient-strict-henselization: sources/stacks-spaces-more-morphisms/pinned-original-source.tex, source lines 5787--5787. Referenced by the retained statement or proof but not part of the collected spans; supplied only as a navigation pointer. The argument does not depend on its content.\par\medskip
\ifdefined\hypertarget\hypertarget{sourcepointer-2}{}\fi
\noindent\textit{Source reference.} Original lemma spaces-morphisms-definition-flat-module: sources/stacks-spaces-more-morphisms/pinned-original-source.tex, source lines 5787--5787. Referenced by the retained statement or proof but not part of the collected spans; supplied only as a navigation pointer. The argument does not depend on its content.\par\medskip
\ifdefined\hypertarget\hypertarget{sourcepointer-3}{}\fi
\noindent\textit{Source reference.} Original lemma spaces-morphisms-lemma-base-change-module-flat: sources/stacks-spaces-more-morphisms/pinned-original-source.tex, source lines 5787--5787. Referenced by the retained statement or proof but not part of the collected spans; supplied only as a navigation pointer. The argument does not depend on its content.\par\medskip
\ifdefined\hypertarget\hypertarget{sourcepointer-4}{}\fi
\noindent\textit{Source reference.} Original lemma spaces-morphisms-lemma-flat-at-point: sources/stacks-spaces-more-morphisms/pinned-original-source.tex, source lines 5787--5787. Referenced by the retained statement or proof but not part of the collected spans; supplied only as a navigation pointer. The argument does not depend on its content.\par\medskip
\ifdefined\hypertarget\hypertarget{sourcepointer-5}{}\fi
\noindent\textit{Source reference.} Original lemma spaces-properties-lemma-describe-etale-local-ring: sources/stacks-spaces-more-morphisms/pinned-original-source.tex, source lines 5787--5787. Referenced by the retained statement or proof but not part of the collected spans; supplied only as a navigation pointer. The argument does not depend on its content.\par\medskip
\ifdefined\hypertarget\hypertarget{sourcepointer-6}{}\fi
\noindent\textit{Source reference.} Original lemma spaces-properties-lemma-geometric-lift-to-usual: sources/stacks-spaces-more-morphisms/pinned-original-source.tex, source lines 5787--5787. Referenced by the retained statement or proof but not part of the collected spans; supplied only as a navigation pointer. The argument does not depend on its content.\par\medskip
\ifdefined\hypertarget\hypertarget{sourcepointer-7}{}\fi
\noindent\textit{Source reference.} Original lemma spaces-properties-lemma-stalk-quasi-coherent: sources/stacks-spaces-more-morphisms/pinned-original-source.tex, source lines 5787--5787. Referenced by the retained statement or proof but not part of the collected spans; supplied only as a navigation pointer. The argument does not depend on its content.\par\medskip
\begin{lemma}
\label{lemma-flat-on-fibres-at-point}
In the situation above the following are equivalent
\begin{enumerate}
\item Pick a geometric point $\overline{x}$ of $X$ lying over $x$.
Set $\overline{y} = f \circ \overline{x}$ and
$\overline{z} = g \circ \overline{x}$. Then the module
$\mathcal{F}_{\overline{x}}/
\mathfrak m_{\overline{z}}\mathcal{F}_{\overline{x}}$
is flat over
$\mathcal{O}_{Y, \overline{y}}/
\mathfrak m_{\overline{z}}\mathcal{O}_{Y, \overline{y}}$.
\item Pick a morphism $x : \Spec(K) \to X$ in the equivalence class of
$x$. Set $z = g \circ x$, $X_z = \Spec(K) \times_{z, Z} X$,
$Y_z = \Spec(K) \times_{z, Z} Y$, and $\mathcal{F}_z$ the pullback
of $\mathcal{F}$ to $X_z$. Then $\mathcal{F}_z$ is flat at $x$ over
$Y_z$ (as defined in Morphisms of Spaces,
Definition \ref{spaces-morphisms-definition-flat-module}).
\item Pick a commutative diagram
$$
\xymatrix{
& & & U \ar[llld]_a \ar[rr] \ar[dr] & & V \ar[llld]_>>>>>>>b \ar[dl] \\
X \ar[rr]_f \ar[dr]_g & & Y \ar[dl]^h &  & W \ar[llld]_c \\
& Z
}
$$
where $U, V, W$ are schemes, and $a, b, c$ are \'etale,
and a point $u \in U$ mapping to $x$. Let $w \in W$ be the image of
$u$. Let $\mathcal{F}_w$ be the pullback of $\mathcal{F}$ to
the fibre $U_w$ of $U \to W$ at $w$. Then $\mathcal{F}_w$
is flat over $V_w$ at $u$.
\end{enumerate}
\end{lemma}

\begin{proof}
Note that in (2) the morphism $x : \Spec(K) \to X$ defines a
$K$-rational point of $X_z$, hence the statement makes sense. Moreover,
the condition in (2) is independent of the choice of $\Spec(K) \to X$
in the equivalence class of $x$ (details omitted; this will also follow
from the arguments below because the other conditions do not depend
on this choice). Also note that we can always choose a diagram as in
(3) by: first choosing
a scheme $W$ and a surjective \'etale morphism $W \to Z$, then choosing
a scheme $V$ and a surjective \'etale morphism $V \to W \times_Z Y$, and
finally choosing a scheme $U$ and a surjective \'etale morphism
$U \to V \times_Y X$. Having made these choices we set $U \to W$ equal
to the composition $U \to V \to W$ and we can pick a point $u \in U$ mapping
to $x$ because the morphism $U \to X$ is surjective.

\medskip\noindent
Suppose given both a diagram as in (3) and a geometric point
$\overline{x} : \Spec(k) \to X$ as in (1). By
Properties of Spaces, Lemma
\ref{spaces-properties-lemma-geometric-lift-to-usual}
we can choose a geometric point $\overline{u} : \Spec(k) \to U$
lying over $u$ such that $\overline{x} = a \circ \overline{u}$.
Denote $\overline{v} : \Spec(k) \to V$ and
$\overline{w} : \Spec(k) \to W$ the induced geometric points of
$V$ and $W$. In this setting we know that
$\mathcal{O}_{X, \overline{x}} = \mathcal{O}_{U, u}^{sh}$
and similarly for $Y$ and $Z$, see
Properties of Spaces,
Lemma \ref{spaces-properties-lemma-describe-etale-local-ring}.
In the same vein we have
$$
\mathcal{F}_{\overline{x}} =
(a^*\mathcal{F})_u \otimes_{\mathcal{O}_{U, u}}
\mathcal{O}_{U, u}^{sh}
$$
see
Properties of Spaces, Lemma \ref{spaces-properties-lemma-stalk-quasi-coherent}.
Note that the stalk of $\mathcal{F}_w$ at $u$ is given by
$$
(\mathcal{F}_w)_u = (a^*\mathcal{F})_u/\mathfrak m_w(a^*\mathcal{F})_u
$$
and the local ring of $V_w$ at $v$ is given by
$$
\mathcal{O}_{V_w, v} = \mathcal{O}_{V, v}/\mathfrak m_w\mathcal{O}_{V, v}.
$$
Since $\mathfrak m_{\overline{z}} =
\mathfrak m_w \mathcal{O}_{Z, \overline{z}} =
\mathfrak m_w \mathcal{O}_{W, w}^{sh}$
we see that
\begin{align*}
\mathcal{F}_{\overline{x}}/
\mathfrak m_{\overline{z}}\mathcal{F}_{\overline{x}} & =
(a^*\mathcal{F})_u \otimes_{\mathcal{O}_{U, u}}
\mathcal{O}_{X, \overline{x}}/
\mathfrak m_{\overline{z}}\mathcal{O}_{X, \overline{x}} \\
& =
(\mathcal{F}_w)_u \otimes_{\mathcal{O}_{U_w, u}}
\mathcal{O}_{U, u}^{sh}/\mathfrak m_w\mathcal{O}_{U, u}^{sh} \\
& = (\mathcal{F}_w)_u \otimes_{\mathcal{O}_{U_w, u}}
\mathcal{O}_{U_w, \overline{u}}^{sh} \\
& = (\mathcal{F}_w)_{\overline{u}}
\end{align*}
the penultimate equality by
Algebra, Lemma \ref{algebra-lemma-quotient-strict-henselization}
and the last equality by
Properties of Spaces, Lemma \ref{spaces-properties-lemma-stalk-quasi-coherent}.
The same arguments applied to the structure sheaves of $V$ and $Y$
show that
$$
\mathcal{O}_{V_w, \overline{v}}^{sh} =
\mathcal{O}_{V, v}^{sh}/\mathfrak m_w \mathcal{O}_{V, v}^{sh} =
\mathcal{O}_{Y, \overline{y}}/
\mathfrak m_{\overline{z}}\mathcal{O}_{Y, \overline{y}}.
$$
OK, and now we can use
Morphisms of Spaces, Lemma \ref{spaces-morphisms-lemma-flat-at-point}
to see that (1) is equivalent to (3).

\medskip\noindent
Finally we prove the equivalence of (2) and (3).
To do this we pick a field extension $\tilde K$ of $K$
and a morphism $\tilde x : \Spec(\tilde K) \to U$ which
lies over $u$ (this is possible because $u \times_{X, x} \Spec(K)$
is a nonempty scheme). Set $\tilde z : \Spec(\tilde K) \to U \to W$
be the composition. We obtain a commutative diagram
$$
\xymatrix{
& & & U_w \times_w \tilde z \ar[llld]_a \ar[rr] \ar[dr] & &
V_w \times_w \tilde z \ar[llld]_>>>>>>>b \ar[dl] \\
X_z \ar[rr]_f \ar[dr]_g & & Y_z \ar[dl]^h &  & \tilde z \ar[llld]_c \\
& z
}
$$
where $z = \Spec(K)$ and $w = \Spec(\kappa(w))$. Now it
is clear that $\mathcal{F}_w$ and $\mathcal{F}_z$ pull back to the
same module on $U_w \times_w \tilde z$. This leads to a commutative
diagram
$$
\xymatrix{
X_z \ar[d] & U_w \times_w \tilde z \ar[l] \ar[d] \ar[r] & U_w \ar[d] \\
Y_z & V_w \times_w \tilde z \ar[l] \ar[r] & V_w
}
$$
both of whose squares are cartesian and whose bottom horizontal
arrows are flat: the lower left horizontal arrow is the composition
of the morphism $Y \times_Z \tilde z \to Y \times_Z z = Y_z$ (base change
of a flat morphism), the \'etale morphism
$V \times_Z \tilde z \to Y \times_Z \tilde z$, and
the \'etale morphism $V \times_W \tilde z \to V \times_Z \tilde z$.
Thus it follows from
Morphisms of Spaces,
Lemma \ref{spaces-morphisms-lemma-base-change-module-flat}
that
$$
\mathcal{F}_z\text{ flat at }x\text{ over }Y_z
\Leftrightarrow
\mathcal{F}|_{U_w \times_w \tilde z}
\text{ flat at }\tilde x\text{ over }V_w \times_w \tilde z
\Leftrightarrow
\mathcal{F}_w\text{ flat at }u\text{ over }V_w
$$
and we win.
\end{proof}

\end{document}
